\paragraph{Combined refinement: the 2572 lower bound for closed tables.}
This independently reviewed addendum strengthens the
preceding intermediate bound. The full argument below
excludes twelve more necessary base rows and proves the
closed-table interval $[2572,2744]$.
Its first step excludes the two
NN-G035 base rows
\[
 (w,g_1,g_2,d_1,d_2,u_1,u_2)
   =(8,72,96,6,6,4,8),\quad(8,80,96,8,4,4,8),
\]
with every allowed auxiliary increment, not only their
least auxiliary masses. The mass-ten rows with first-cell
triple $(80,8,4)$ are not excluded by this argument.

In either row, $U_1=D$ and the earlier integral quotient
argument gives $G_1=DH$, $G_2=DWC$,
$H=\delta_1+DC$, $C\equiv D-1\pmod2$, and $C(4)=3$.
The mass-eight $W$ is irreducible and coprime to $D$.
It cannot divide $H$, whose mass is respectively $18$
or $20$, because eight divides neither number.
Thus $W$ is coprime to $G_1$ and both second-cell
identities give $U_2=WR,V_2=WZ$ integrally.
Their parities give $R\equiv Z\equiv1$.
The equal-mass lemma forces $U_2=W$, and hence
\[
             \delta_1+\delta_2=(W-D)C.               \tag{G036.R1}
\]

For the first row, both cells have their unique
mass-six character $\delta_i=\psi+2$, where
$\psi=D(D-3)$. If $W=D+4$, (G036.R1) would give
$C=(\psi+2)/2$, which is not an integral polynomial.
Otherwise $W=D+2T_j$ and
$T_jC=\psi+2=(D-1)(D-2)$.
Since $C\equiv D-1$ has positive degree, this forces
$j=1$, $C=D-1$, $W=3D-4$ and $H=2(D-1)^2$.
Now
\[
                 A_2=D^2WC-DHZ.
\]
If $Z$ has degree one, its leading coefficient is
positive even, so the leading coefficient of $A_2$
is at most $3-4<0$. Higher degrees also fail.
Thus $Z$ is constant, but then $V_2=WZ$ has degree
one and cannot dominate its required $q^2$ coefficient
two. This excludes the first row.

For the second row, $\delta_2=\psi$ and
$F=(\delta_1+\delta_2)/2$ is effective, has mass six,
and dominates $\psi$. The mass-eight part of the
small-mass argument in NN-G036 forces $C=D-1$;
the only possible $W$ are $3D-4$ and $D+2T_2$.
For $W=3D-4$ one has $\delta_1=\psi+4$ and
$H=2D(D-2)+4$. Its degree and leading coefficient
give the same contradiction for $Z$ as in the first row.

For $W=D+2T_2$, the forced polynomials are
\[
 \delta_1=\psi+2T_3,\qquad
 H=2D(D-2)+2T_3,\qquad C=D-1.
\]
Here $D^2WC$ has degree five and leading coefficient
two, while $DH$ has degree four and leading coefficient
two. A nonconstant $Z\equiv1\pmod2$ would again give
a negative leading coefficient in $A_2$.
Consequently $Z$ is a positive odd integer.
The scalar relations give $v_2=8Z$ and
$a_2=384-10v_2$. Since $v_2\ge16$ and $a_2\ge16$,
the only possible odd value is $Z=3$, so
$v_2=24,a_2=144$.

The first auxiliary relation is $a_1=320-12v_1$.
It gives $v_1\le26$; and substitution of every other
now-fixed mass into the total gives
\[
                    T=9760+148(26-v_1).
\]
Since $v_1$ is even, the cap $T\le9912$ forces
$v_1=26$ and $a_1=8$.
We show that these masses cannot satisfy coefficientwise
domination.

Effectivity of $A_1=D^2H-D(D-1)V_1$ gives
$\deg V_1\le3$. Write $V_1=(d,a,b,c)$ in symmetric
Laurent notation. Domination and parity require
$d\ge6$ even, $a\ge5$ odd, $b\ge0$ even, and
$c\ge0$ even. The leading coefficient of $A_1$
also gives $c\le2$.
The mass condition is $d+2a+2b+2c=26$.
The relevant coefficients of $A_1$, after eliminating
$d$, are
\[
 \begin{array}{c|l}
 q^0&-44+2a+6b+8c,\\
 q^1&-24+a+3b+5c,\\
 q^2&14-a-2b-c,\\
 q^3&24-a-3b-4c.
 \end{array}
\]
The target lower bounds at degrees zero and three,
respectively two and one, force
$a+3b+4c=23$. The degree-one target then gives
$c-1\ge1$, so $c=2$ and $a+3b=15$.
The degree-two target gives $b-3\ge1$, hence $b\ge4$.
This implies $a\le3$, contradicting $a\ge5$.
The second row is therefore impossible as well.
\hfill$\square$

Deleting these two rows alone leaves $36$ base rows,
$311$ auxiliary increment pairs, and $19$ first-cell
mass triples, with least relaxed score $1956$.
The following uniform arguments exclude three more
mass blocks.

\paragraph{A common-factor observation.}
The identities used throughout are
\[
 WS=G_1G_2,\qquad U_1G_2=W(G_1-D\delta_1),\qquad
 U_2G_1=W(G_2-D\delta_2).                            \tag{G036.R2}
\]
If a primitive irreducible factor $P$ of $W$ divides
neither $U_1$ nor $U_2$, then it divides both $G_1$
and $G_2$. Indeed it divides at least one from the
first identity, and the appropriate one of the next
two forces it into the other. All divisibilities
first obtained over $\mathbb Q[D]$ are integral by
primitivity. In particular divisibility by $P$ forces
$P(4)$ to divide the corresponding ordinary mass.

\paragraph{The entire mass block $(w,g_1,g_2)=(10,80,80)$.}
Here $u_i=10-d_i/2$, so $u_i\in\{4,6,8\}$.
If $W$ is irreducible, its mass ten divides neither
$u_i$, and (G036.R2) gives $W\mid G_1,G_2$.
If $W=(D-2)B$ is the exceptional case of (G036.2),
the primitive irreducible $B$ has mass five, again
dividing neither $u_i$. Thus write $G_i=BP_i$.
The second and third identities give
\[
 U_1P_2=(D-2)(BP_1-D\delta_1),\qquad
 U_2P_1=(D-2)(BP_2-D\delta_2).
\]
An effective character of mass four, six or eight,
dominating and congruent to $D$, has nonzero value
at $D=2$: these values are $2$, $4$, or one of
$6,2,-2$, respectively. Therefore $D-2$ divides
both $P_i$. We have proved $W\mid G_1,G_2$ also
in this exceptional case.

Now write $G_i=WP_i$ in every case. Then
\[
 P_1(4)=P_2(4)=8,\qquad P_1\equiv0,\quad
 P_2\equiv D(D-1)\pmod2.
\]
The polynomial $P_2$ cannot be divisible by $D$:
its integral quotient would have value two at $D=4$
but would reduce to $D-1$ modulo two, whose value
there is odd. Thus $P_2(0)\ne0$.
The divided identities at $D=0$ show also
$P_1(0)\ne0$, because $W(0)\ne0$.
Multiplying them gives
\[
                        W(0)^2=U_1(0)U_2(0).        \tag{G036.R3}
\]
For the mass-ten characters, $W(0)$ is six or
minus two. For the possible $U_i$, the values are
zero, two, four or minus four. Equality (G036.R3)
is possible only when $W(0)=-2$ and both $U_i(0)=2$.
Consequently
\[
 U_1=U_2=D+2,\qquad d_1=d_2=8,\qquad
 W=D+2+2T_j\quad\hbox{with }j\hbox{ odd}.             \tag{G036.R4}
\]
Each cell character is then $\delta_i=\psi+2R_i$,
where $R_i$ has mass two and is either the constant
two or some $T_k$.

Subtracting the two divided identities in (G036.R2),
and putting $E=R_1-R_2$, gives
\[
 A_0(P_1-P_2)=E,\qquad
 A_0=\frac{D+2+T_j}{D}.                             \tag{G036.R5}
\]
The quotient is integral since $j$ is odd and
$T_j(0)=-2$. The difference $P_1-P_2$ is nonzero,
as its reduction is $D(D-1)$, so $E\ne0$.
Every nonzero difference between two members of
$\{2,T_1,T_2,\ldots\}$ has leading coefficient
$1$ or $-1$, and its value at $D=0$ has absolute
value at most four.

If $j=1$, $A_0=2$, contradicting this odd leading
coefficient of $E$.
For odd $j\ge3$, $A_0$ is monic of degree $j-1$.
Thus $P_1-P_2$ has leading coefficient $\pm1$;
its reduction $D(D-1)$ forces its actual degree
to be exactly two. Write it $aD^2+bD+c$, with
$a=\pm1$, $b$ odd and $c$ even.
The recurrence $T_{j+1}=(D-2)T_j-T_{j-1}$ gives
\[
 T_j'(0)=(-1)^{j-1}j^2,\qquad A_0(0)=1+j^2.
\]
For completeness, the derivative recurrence is
$r_{j+1}=2(-1)^j-2r_j-r_{j-1}$ with
$r_0=0,r_1=1$; the displayed formula satisfies it
by direct substitution.
Evaluating (G036.R5) at zero now gives
$|(1+j^2)c|\le4$, so $c=0$.
But $P_1(4)=P_2(4)$ then implies
$16a+4b=0$, making $b=-4a$ even, a contradiction.
This excludes the entire mass block, including
all six base rows and all their auxiliary choices.

\paragraph{The entire mass block $(w,g_1,g_2)=(8,96,96)$.}
Here $u_i=8-d_i/3$ is four or six.
The mass-eight $W$ is irreducible, and its mass
divides neither $u_i$, so (G036.R2) again gives
$G_i=WP_i$. Now $P_i(4)=12$, with
$P_1\equiv0$ and $P_2\equiv D(D-1)$ modulo two.
The polynomial $P_1$ cannot be divisible by $D$:
the quotient would have even coefficients but
ordinary mass $12/4=3$.
The divided identities at zero consequently
force both $P_i(0)\ne0$ and give (G036.R3).
But $W(0)^2=16$, whereas $U_i(0)$ are zero or two,
whose products are zero or four. This contradiction
excludes both base rows in this block.

\paragraph{The two asymmetric mass-ten blocks.}
The finite certificate has one base row for
$(w,g_1,g_2)=(10,76,100)$, with
$(d_1,d_2,u_1,u_2)=(4,6,6,10)$, and its index-swapped
row with $(g_1,g_2)=(100,76)$.
Let $G_-$ denote the smaller action character and
$G_+$ the larger. Their cell characters are
respectively $\psi$ and $\psi+2$, and the smaller
index has $U_-=D+2$.

If $W$ is irreducible, it cannot divide $G_-$,
since ten does not divide $76$.
It therefore divides $G_+$ and then $U_+$ by
(G036.R2). Since $u_+=10$, the equal-mass lemma
forces $U_+=W$.
If $W=(D-2)B$, the irreducible mass-five $B$ cannot
divide $G_-$, so the same argument gives $B\mid U_+$.
Writing $U_+=BR$, its reduction is $R\equiv D$.
The leading-coefficient argument from NN-G036 forces
$\deg R=1$ and then the mass-ten classification gives
$U_+=W$ in this case as well.

The larger-index identity now gives
$G_+=G_-+D(\psi+2)$.
Substitution into the smaller-index identity yields
\[
 (W-D-2)G_-=D\bigl(W\psi+(D+2)(\psi+2)\bigr).       \tag{G036.R6}
\]
If $W=D+6$, its denominator is four and the resulting
$G_-$ has leading coefficient $1/2$, contradicting
integrality.
For $W=D+2+2T_j$, equation (G036.R6) gives
\[
 G_-=D\psi+\frac{N(D)}{T_j(D)},\qquad
 N(D)=D(D+2)(D^2-3D+1).
\]
Integrality forces the monic $T_j$ to divide $N$.
Thus $j\le4$. None does:
$N(2)=-8$ excludes $T_1=D-2$ and also $T_3$,
which contains $T_1$;
the remainder modulo $T_2=D^2-4D+2$ is $16D-10$;
and $T_4=D^4-8D^3+20D^2-16D+2$ is not the
monic degree-four polynomial $N$.
This excludes both asymmetric blocks, uniformly
over all original support indices.
\hfill$\square$

\paragraph{The resulting finite dimension gate.}
In addition to the five rows already removed in
NN-G036, this addendum removes two isolated mass-eight
rows, six rows in the $(10,80,80)$ block, two in the
$(8,96,96)$ block, and the two asymmetric rows.
The remaining complete necessary projection has
$26$ base rows, $42$ auxiliary increment pairs,
and $15$ first-cell mass triples. Its least relaxed
score is $2572$. The proved exact attainer still
has score $2744$, so the closed-table minimum lies
in $[2572,2744]$.

In the NN-G010 geometry, these are exclusions of
ordinary-dimension patterns for actual composing
graded sign-representation objects. The tests at
$D=0$ are integral character evaluations at $q=-1$;
the derivative there is a character-polynomial
calculation, not a Weil-weight assertion.
They preserve the SGA 3, I, 4.7.2--4.7.3 and
archive cote 2, scan 17 dictionary of NN-G036.
The surviving rows are not declared realizable,
and this lower bound is not transferred to arbitrary
extra-position lifts or canonical GCT multiplicities.

